\section{Scope, conventions, and the constructive reading}\label{sec:scope}

\subsection{The three source pages}

The requested source material consists of the Wikipedia articles
\emph{Bring radical} and \emph{Quintic equation}, and the MathWorld article
\emph{Quintic Equation}.  They mix several logically different kinds of
assertion:

\begin{enumerate}
\item finite algebraic identities, such as a Tschirnhaus substitution, a
      resultant, a cyclotomic product, or a displayed radical formula;
\item existence and approximation assertions for real or complex roots;
\item convergence, branch-point, differential-equation, theta-function, and
      hypergeometric assertions;
\item Galois-theoretic impossibility and solvability criteria;
\item historical attributions and computational advice.
\end{enumerate}

Only the first four are mathematical propositions.  Historical priority is not
an arithmetic theorem and is therefore recorded only as bibliography.  The
three pages were inspected on 16 November 2026.  Their main claims relevant to
this report are the following.

\begin{itemize}
\item The real Bring radical \(\BR(a)\) is the unique real zero of
      \(x^5+x+a\); it is odd, decreasing and unbounded, with
      \(\BR(a)\sim-a^{1/5}\).  Its local series, hypergeometric form, radius,
      and four branch points are displayed.
\item A general quintic can first be depressed, then put in principal form,
      and finally related to a Bring--Jerrard equation by a quartic
      Tschirnhaus transformation.  A trinomial \(x^5+px+q\) is reduced to one
      variable of the Bring radical by scaling.
\item General quintics are not solvable by radicals, while special quintics
      are.  The pages give de Moivre's family, resolvent constructions,
      solvable parameterizations, explicit examples, and finite-group
      criteria.
\item Beyond radicals, the pages describe hypergeometric, theta, elliptic,
      icosahedral and iterative methods, and the quintic arising from the
      collinear Lagrange points.
\end{itemize}

The goal is not to copy those pages, but to identify which claims admit a clean
finitary proof and to expose the exact data needed by \ipra.

\subsection{Four meanings of ``proved in \(\ipra\)''}

First-order \ipra{} quantifies only over natural numbers.  It has no primitive
sort of arbitrary real numbers, complex functions, fields, or infinite
series.  We therefore use four precise readings.

\begin{definition}[Exact finite statement]
An \emph{exact} statement is an identity or decidable relation on natural-number
codes for integers, rationals, finite vectors, polynomials, permutations, or
finite tables.  Denominators are cleared and all loops have explicit bounds.
Such a statement is directly expressible in \ipra.
\end{definition}

\begin{definition}[Effective analytic statement]
A real or complex quantity is represented by a fixed primitive-recursive
sequence of rational or Gaussian-rational approximants, together with a fixed
primitive-recursive modulus.  An analytic theorem is a theorem schema for such
named data, with all radii, bounds and error moduli explicit.  No quantifier
ranges over an unspecified completed function space.
\end{definition}

\begin{definition}[Certificate statement]
A certificate statement says that a primitive-recursive verifier accepts a
finite object only when a specified mathematical relation holds.  Examples are
a resultant identity, a factorization modulo a prime, a multiplication table
for a finite Galois group, or a Tschirnhaus transformation.  The theorem is
uniform in the certificate code; a concrete certificate is then checked by
finite computation.
\end{definition}

\begin{definition}[Formalization outline]
An outline gives all object-theory definitions, the exact lemmas to be proved,
and a dependency path to a kernel-checked Lean theorem, but does not claim that
the final Lean term has already been compiled.  We use this label for the
large global theta-function and universal Bring--Jerrard developments.
\end{definition}

The audit in Section~\ref{sec:audit} labels results by
\(\Exact\), \(\Effective\), \(\Conditional\), or \(\Outline\).  These labels
are deliberately stronger than a vague assertion that an ordinary classical
proof is ``finitary in spirit.''

\subsection{Sign conventions}

Two closely related normal forms occur in the sources.  Throughout this report
we define
\begin{equation}\label{eq:BRdef}
  \BR(a)^5+\BR(a)+a=0
\end{equation}
for real \(a\).  Thus if
\[
  h(x)=x^5+x,
\]
then \(\BR(a)=h^{-1}(-a)=-h^{-1}(a)\).  The alternative Bring equation
\(t^5-t+\rho=0\) is treated separately; it has different critical values and
is the convention used in the Cockle--Harley differential resolvent.

For a primitive fifth root of unity we write \(\omega\), so
\begin{equation}\label{eq:cyclo5}
 \omega^5=1,
 \qquad 1+\omega+\omega^2+\omega^3+\omega^4=0.
\end{equation}
No analytic choice of \(e^{2\pi i/5}\) is needed for the algebraic identities:
\(\omega\) is simply an element satisfying these two equations.

\subsection{What is new in this packet}

Several ingredients already existed in ProveIt.  The repository's
\code{IPRADerivations} library contains kernel-checked object-theory proofs of
the quadratic, Cardano and Ferrari factorization identities.  Its Lagrange
inversion packet gives the finite coefficient algebra needed for inverses of
\(x-x^p\).  A separate polynomial-formulas packet develops a generic
Abel--Ruffini architecture.  We use these rather than silently reproving them
in a stronger metatheory.

The new contribution is the synthesis specific to Bring radicals and the
quintic pages:

\begin{itemize}
\item a short sum-of-squares proof of strong monotonicity and the resulting
      primitive-recursive real Bring radical;
\item exact coefficient and branch-point proofs tied to the finite Lagrange
      framework;
\item a transparent principal-form reduction and a compact certificate
      interface for the large Bring--Jerrard step;
\item concrete finite-field certificates proving that both
      \(x^5-x+1\) and \(x^5+x+3\) have Galois group \(S_5\);
\item exact de Moivre and four-resolvent identities, plus a complete finite
      enumeration of the subgroup classification relevant to solvable
      irreducible quintics;
\item a detailed Lean module plan connected to the existing ProveIt theorem
      names.
\end{itemize}

The executable verifier performs seventeen independent families of exact
checks.  It is evidence for the displayed certificates, not a substitute for
the human proofs or the intended object-theory derivations.
